\scanpage{24}
\begin{center}-- 34 --\end{center}
\hfill n° 371
\[
\begin{aligned}
w((\bigotimes_i x_i)(\bigotimes_i y_i))&=w(\bigotimes_i(x_i y_i))=\prod_i v_i(x_i y_i)\\
&=\prod_i(v_i(x_i)v_i(y_i))=(\prod_i v_i(x_i))(\prod_i v_i(y_i))
\end{aligned}
\]
car les $v_i(y_i)$ permutent deux à deux et chaque $v_i(y_i)$ permute avec tous les $v_j(x_j)$ tels que $j\ne i$. Il est clair que l'on a $v_i=w\circ u_i$ pour tout $i\in I$, et l'unicité de $w$ résulte de ce que les $u_i(E_i)$ engendrent $E$.

On dit que les homomorphismes $u_j$ sont \uline{canoniques}, et on écrit $w=\bigotimes_i v_i$ s'il n'en résulte pas de confusion.

Le couple formé de $E$ et de l'application canonique $\varphi:(x_i)\mapsto\bigotimes_i x_i$ de $\prod_i E_i$ dans $E$ est solution du \uline{problème d'application universelle} (\uline{Ens.}, chap. IV, § 3, n° 1) où $\Sigma$ est l'espèce de structure de $A$-algèbre, les morphismes étant les $A$-homomorphismes d'algèbres, et les $\alpha$-applications les applications $\prod_i u_i$ de $\prod_i E_i$ dans une $A$-algèbre, telles que les $u_i$ soient des homomorphismes d'algèbres et que $u_i(x_i)$ et $u_j(x_j)$ permutent pour $i\ne j$, quels que soient $x_i\in E_i$ et $x_j\in E_j$.

\uline{COROLLAIRE.}-- \uline{Soient} $(E_i)_{i\in I}$, $(F_i)_{i\in I}$ \uline{deux familles finies de} $A$-\uline{algèbres}, \uline{et pour tout} $i\in I$, \uline{soit} $f_i:E_i\to F_i$ \uline{un homomorphisme d'algèbres}. \uline{Si} $u_i:E_i\to\bigotimes_i E_i$, $v_i:F_i\to\bigotimes_i F_i$ \uline{sont les homomorphismes canoniques}, \uline{l'application} $f=\bigotimes_i f_i$ \uline{est l'unique homomorphisme de} $A$-\uline{algèbres tel que} $f\circ u_i=v_i\circ f_i$ \uline{pour tout} $i\in I$.

Il suffit de noter que les homomorphismes $g_i=v_i\circ f_i$ sont tels que $g_i(x_i)$ et $g_j(x_j)$ permutent pour $i\ne j$, $x_i\in E_i$ et $x_j\in E_j$; on peut alors appliquer à ces homomorphismes la prop. 5.

Lorsque les $E_i$ sont \uline{commutatives}, il en est de même de $E$, et les hypothèses sur les $v_i$ dans la prop. 5, (ii) sont automatiquement vérifiées si $F$ est elle aussi une $A$-algèbre \uline{commutative}.

\begin{quote}
On peut donc dire que dans ce cas $(E,\varphi)$ est aussi solution du problème d'application universelle où $\Sigma$ est l'espèce de structure de $A$-algèbre \uline{commutative}, les morphismes étant les $A$-homomorphismes
\end{quote}
\mspage{23}
